% open-logic.tex
% compiles to produce complete text using open-logic-dev style

\documentclass[../include/open-logic-part]{subfiles}

\begin{document}

\clearpage

\begin{editorial}
دا دوتنه ټول هغه مطالب رااخلي چې د خلاص منطق په پروژه کښې شامل دي۔ که دا ډول ادارتي يادښتونه ښکاري، نو دا څرګندوي چې دوتنه د دې مطالبو د وړاندې کولو طريقې ته له پام کولو پرته جوړه شوې ده۔ که دا لولئ، نو غالباً له دې پي ډي اېف څخه درس ورکول يا مطالعه کول \emph{مناسب نۀ دي}۔

د خلاص منطق پروژه ډېرې داسې طريقې وړاندې کوي چې د هغو په وسيله د درس ورکولو يا ځان زده کړې دپاره مناسب متن جوړېدل ممکن وي۔ د بېلګې په توګه، په اصلي تنظيم کښې متن ټول منطقي عملګر بنسټيز ګڼي، او په ثبوتونو کښې د ټولو عملګرو ټول حالتونه څېړي۔ خو دا ډېر غوره دي چې ځينې حالتونه د تمرينونو په توګه پرېښودل شي۔ د خلاص منطق پروژه پخپله هم روان کار دے۔ د ګډ کار او ښه والي د هڅولو دپاره، مطالب آن هغه وخت هم شاملېږي چې يوازې د مسودې په بڼه وي، تمرينونه نۀ لري، او داسې نور۔ هغه پي ډي اېف چې د يو کورس دپاره جوړېږي، دا برخې به ترې وباسي۔

د درس او مطالعې دپاره د لا مناسبو پي ډي اېفونو د موندلو دپاره، \href{http://builds.openlogicproject.org/}{د خلاص منطق پروژې په وېبپاڼه شته نمونوي کورسونه} وګورئ۔ د خپل متن د جوړولو دپاره، له \href{https://github.com/OpenLogicProject/OpenLogic/tree/master/courses/sample}{نمونوي چلوونکې دوتنې} څخه پيل کولے شئ، يا د لا پرمختللو او مفصلو بېلګو دپاره له دې پروژې څخه د جوړو شوو درسي کتابونو سرچينې کتلے شئ۔
\end{editorial}

\olimport[sets-functions-relations]{sets-functions-relations-complete}

\olimport[propositional-logic]{propositional-logic}

\olimport[first-order-logic]{first-order-logic}

\olimport[model-theory]{model-theory}

\olimport[computability]{computability}

\olimport[turing-machines]{turing-machines}

\olimport[incompleteness]{incompleteness}

\olimport[second-order-logic]{second-order-logic}

\olimport[lambda-calculus]{lambda-calculus}

\olimport[many-valued-logic]{many-valued-logic}

\olimport[normal-modal-logic]{normal-modal-logic}

\olimport[applied-modal-logic]{applied-modal-logic}

\olimport[intuitionistic-logic]{intuitionistic-logic}

\olimport[counterfactuals]{counterfactuals}

\olimport[set-theory]{set-theory}

\olimport[methods]{methods}

\olimport[history]{history}

\olimport[reference]{reference}

\end{document}
